\section{Introduction: the questions a clinician asks at the bedside}
\label{sec:intro}

This framework is written for a reader who treats cancer patients and who will be
asked, soon, to supervise an autonomous system in a clinical trial, but who has
never driven a surgical robot or written a line of verification code. It defines
its terms as it goes, and it organizes the decision to trust around eight plain
questions a clinician already asks of any new tool.

A few definitions make the rest readable. \textit{Physical AI} is artificial
intelligence that acts in the physical world through a machine, such as a surgical
humanoid in an oncology trial, rather than only returning text on a screen.
\textit{Verification before generation} is the principle that a machine may
propose an action, but the proposal is checked against fixed, published standards
before anything reaches a patient. \textit{Claude Code} is the software agent that
drafts the candidate action, \textit{Codex} is a second, independent agent that
reviews it, and \textit{VVUQ}, meaning verification, validation, and uncertainty
quantification, is the ten-gate examination the proposal must pass. The
examination resolves in exactly one of three ways: ACCEPT and act under audit,
ESCALATE to a qualified human, or BLOCK before execution. \textit{Calibrated
trust} is the goal of this framework: reliance proportioned to demonstrated
capability, avoiding both \textit{automation bias} (over-trust) and
\textit{algorithm aversion} (under-trust).

[DRAFTING INSTRUCTIONS: in the full-framework, expand these definitions with the
mechanism described in the prior review (\texttt{review/}) and the Bill
(\texttt{Kawchak2026HR9510Bill}, DOI 10.5281/zenodo.20619762). Anchor calibrated
trust, automation bias, and algorithm aversion in
\texttt{adoption/references/framework\_refs.bib}
(\texttt{lee2004trust}, \texttt{goddard2012automation},
\texttt{dietvorst2015algorithm}). Reproduce Mermaid figures
\texttt{adoption/mermaid/diagrams/03-verification-before-generation.md} and
\texttt{02-eight-confidence-questions.md} as colored \texttt{cfwfig} TikZ.]

\begin{cfwfig}
\node[cfone] (a) at (0,0) {Clinical question};
\node[cftwo] (b) at (3.4,0) {Claude Code proposes};
\node[cftwo] (c) at (6.8,0) {Codex reviews};
\node[cfdec] (d) at (10.0,0) {Ten gates};
\node[cfhope] (acc) at (5.0,-2.4) {ACCEPT under audit};
\node[cfthree] (esc) at (8.4,-2.4) {ESCALATE};
\node[cfrisk] (blk) at (11.4,-2.4) {BLOCK};
\draw[cfedge] (a)--(b); \draw[cfedge] (b)--(c); \draw[cfedge] (c)--(d);
\draw[cfedge] (d) to[out=-120,in=60] (acc.north);
\draw[cfedge] (d) -- (esc.north);
\draw[cfedge] (d) to[out=-60,in=120] (blk.north);
\end{cfwfig}
\vspace{-0.15cm}
\cfwcaption{Figure 1. Verification before generation, the clinician's view
(placeholder; expanded from Mermaid 03 in the full-framework).}

\begin{cfwfig}
\node[cfone,text width=15mm]  (q1) at (0,0)    {1 Comp};
\node[cfrisk,text width=15mm] (q2) at (2.7,0)  {2 Safe};
\node[cftwo,text width=15mm]  (q3) at (5.4,0)  {3 Trans};
\node[cftwo,text width=15mm]  (q4) at (8.1,0)  {4 Over};
\node[cfrisk,text width=15mm] (q5) at (8.1,-1.9){5 Equity};
\node[cfthree,text width=15mm](q6) at (5.4,-1.9){6 Rely};
\node[cftwo,text width=15mm]  (q7) at (2.7,-1.9){7 Acct};
\node[cfhope,text width=15mm] (q8) at (0,-1.9) {8 Flow};
\node[cfact,text width=20mm]  (ct) at (4.0,-3.5){Calibrated trust};
\draw[cfedge] (q1)--(q2); \draw[cfedge] (q2)--(q3); \draw[cfedge] (q3)--(q4);
\draw[cfedge] (q4)--(q5); \draw[cfedge] (q5)--(q6); \draw[cfedge] (q6)--(q7);
\draw[cfedge] (q7)--(q8); \draw[cfedge] (q8) to[out=-90,in=180] (ct.west);
\end{cfwfig}
\vspace{-0.15cm}
\cfwcaption{Figure 2. The eight confidence questions (placeholder; expanded from
Mermaid 02 in the full-framework). They accumulate into calibrated trust.}

Three questions recur throughout, asked politely and answered concretely. Would I
let this system near my own family member? Can I defend this decision at tumor
board? And if an auditor asked me tomorrow, could I show exactly what happened and
why? The remainder takes the eight confidence questions in turn, each defining its
terms, pairing its claim with a cited fact, and closing on the specific mechanism
that answers it.
